%% kapa: exam=1 round=20 year=2009 period=1 form=A subject=civil_law subject_ko=민법 seq=3 number=3
%% source: Q-Net 감정평가사 기출문제 (cst00309) artlSeq=5024178 file=extracted/1008/5024178/2009년 제20회 감정평가사 1차시험 기출문제/1교시.hwp
%% conversion: kordoc 4.24.51 -> tools/gen1.py (auto); review=proofread-sonnet; flags=none
\begin{question}{3}
신의성실의 원칙에 관한 설명으로 옳지 \uline{않은} 것은? (다툼이 있으면 판례에 의함)
\begin{choices}
  \item 신의성실의 원칙은 채권관계뿐만 아니라 물권관계나 가족관계에서도 적용된다.
  \item 신의성실의 원칙은 구체적 내용이 정하여져 있지 않은 일반조항으로서 그 내용은 재판에 의하여 형성된다.
  \item 신의성실의 원칙은 강행법규의 성질을 가지므로 당사자의 주장이 없더라도 법원이 직권으로 판단할 수 있다.
  \item 강행법규 위반사실을 알면서 그러한 계약을 체결한 당사자가 후에 그 행위가 강행법규 위반으로 무효라고 주장하는 것은 신의칙에 위반된다.
  \item 신의성실의 원칙에 위배됨을 이유로 권리행사를 부정하기 위해서는 법률관계의 당사자가 상대방의 신의에 반하여 권리를 행사하는 것이 정의관념에 비추어 용인될 수 없는 정도의 상태에 이르러야 한다.
\end{choices}
\end{question}
\answer{4}
